\documentclass[12pt]{article}

\title{Self-Motions of Griffis-Duffy Type Parallel Manipulators}

\author{ Manfred HUSTY$^1$ \& Adolf KARGER$^2$}

\date{\small $^1$Inst. f. Mathematics and Applied Geometry, Montan
  University, Leoben, Austria, {\tt
    e-mail:husty\symbol{64}unileoben.ac.at},\\ $^2$Mathematical
  Institute of the Charles University, Praha, Czech Republic, {\tt
    e-mail:karger\symbol{64}karlin.mff.cuni.cz} }

\begin{document}
\maketitle

Griffis and Duffy proposed in several papers and an US patent
\cite{du} a new type of parallel manipulator. This manipulator
(GD-platform) is a special case of the well known Stewart-Gough
platform which consists of a base and a platform connected by six
ball-joint-ended linear actuators. On base and platform the centers of
the ball joints form two sets of attachment points. Giffis and Duffy
propose a very special geometry for these two sets of points. On both,
base and platform, three points form a triangle and the remaining
three points are placed each on one side of the triangle. The
connection between base and platform is made in such a way that the
points on the sides of the triangle of the base are connected to the
vertices of the triangle of the platform and vice versa.

Griffis and  Duffy have shown that this type of platform has a
relatively simple direct kinematics which admits up to 16 real
soutions.  A closer inspection of this manipulator reveals that it has
some strange and maybe dangerous behaviour. Using an improved version
of the direct kinematics algorithm developed in \cite{hu1} and
\cite{ka1} we will show in this paper

\begin{enumerate}
\item One relative simple design condition will make the GD-platform
  architectural singular (this property says that the manipulator is
  singular in every point of its workspace).
\item The design condition has nice  a geometric interpretation.
\item The GD-platform must obey an assembly condition (if it does
 not obey this assembly condition then there is no real assembly mode).
\item If the platform obeys the assembly condition than it has a self
  motion, which means that it is mobile at least with one degree of
  freedom without changing the actuator lengths.
\end{enumerate}

\begin{thebibliography}{99}


\bibitem{du}GRIFFIS, M. - DUFFY, J.: United States Patent \# 5,179,525
  (1993).

\bibitem{ka1}HUSTY, M. - KARGER, A. - SACHS, H. - STEINHILPER, W.:
  Kinematik und Robotik, Springer Verlag, Berlin - Heidelberg - New
  York, ISBN 3-540-63181-X, 633~S, (1997).

\bibitem{hu1} HUSTY, M. L.: An Algorithm for Solving the Direct
  Kinematics of General Stewart-Gough Platforms. Mech. Mach. Theory
  31, 4 (1996), 365-380.


\bibitem{du2}GRIFFIS, M. - CRANE, C. - DUFFY, J.: A Smart Kinestatic
Interactive Platform, in: J. Lenarcic and B. Ravani (eds.) Advances in
Robot Kinematics and Computational Geometry, Kluwer Acad. Pub. , ISBN
0-7923-2983, (1994) 459-464.
\end{thebibliography}

\end{document}
